\documentclass[11pt,a4paper]{article}

\usepackage[T1]{fontenc}
\usepackage[utf8]{inputenc}
\usepackage{lmodern}
\usepackage[ngerman]{babel}
\usepackage[a4paper,margin=2.2cm]{geometry}
\usepackage{array}
\usepackage{longtable}
\usepackage{tabularx}
\usepackage{xcolor}
\usepackage{hyperref}

\setlength{\parindent}{0pt}
\setlength{\parskip}{0.5em}
\renewcommand{\arraystretch}{1.35}

\newcommand{\answerline}[1]{\rule{#1}{0.4pt}}
\newcommand{\term}[1]{\fcolorbox{black!20}{blue!6}{\strut #1}}

\title{\textbf{Aufgabenblatt: Systemarchitekturen zur Softwareentwicklung}}
\author{Modul 321 -- Lektion 3}
\date{}

\begin{document}

\maketitle

\textbf{Name:} \answerline{7cm} \hfill \textbf{Datum:} \answerline{3.5cm}

\vspace{0.8em}

\fcolorbox{blue!40}{blue!6}{
  \parbox{\dimexpr\linewidth-2\fboxsep-2\fboxrule-6pt\relax}{
    \textbf{Ziel der Lektion:} Sie können zentrale Systemarchitekturen
    unterscheiden, mit passenden Beispielen erklären und Architekturentscheide
    im Spannungsfeld von Komplexität, Skalierung, Betrieb und Teamstruktur begründen.
  }
}

\section*{Hinweise}
\begin{itemize}
  \item Antworten Sie in ganzen Sätzen oder präzisen Stichworten.
  \item Begründen Sie Architekturentscheide knapp, aber fachlich sauber.
  \item Skizzen und kleine Diagramme sind ausdrücklich erlaubt.
\end{itemize}

\section*{1. Begriffe erklären}
Erklären Sie die folgenden Begriffe in \textbf{1 bis 3 eigenen Sätzen}.

\begin{longtable}{|>{\raggedright\arraybackslash}p{0.28\linewidth}|>{\raggedright\arraybackslash}p{0.64\linewidth}|}
\hline
\textbf{Begriff} & \textbf{Ihre Erklärung} \\
\hline
\term{Monolith} & \\[2.2cm] \hline
\term{Microservices} & \\[2.2cm] \hline
\term{Client-Server} & \\[2.2cm] \hline
\term{Peer-to-Peer} & \\[2.2cm] \hline
\end{longtable}

\section*{2. Beispiele zuordnen}
Ordnen Sie jedem Szenario die \textbf{passendste Architektur} zu:

\term{Monolith}, \term{Microservices}, \term{Client-Server}, \term{Peer-to-Peer}

\begin{enumerate}
  \item Eine Webanwendung läuft als ein gemeinsames Deployment mit einer zentralen Datenbank.
  \item Eine Browser-App ruft eine zentrale API auf und erhält dort Daten und Regeln.
  \item Mehrere eigenständige Dienste für Bestellung, Zahlung, Versand und Benachrichtigung kommunizieren über APIs.
  \item Teilnehmer laden Dateien direkt untereinander aus dem Netzwerk herunter.
  \item Eine kleine Schulsoftware wird von einem Team entwickelt und als eine zusammenhängende Anwendung betrieben.
  \item Ein Blockchain-Netzwerk verteilt Daten und Validierung auf viele gleichberechtigte Knoten.
\end{enumerate}

\vspace{0.8em}
\textbf{Begründen Sie danach zwei Ihrer Zuordnungen kurz in je 1 Satz.}

\section*{3. Monolith oder Microservices?}
Ein junges Startup baut eine erste Version einer Terminplattform.

Rahmenbedingungen:
\begin{itemize}
  \item 4 Entwicklerinnen und Entwickler
  \item Anforderungen ändern sich häufig
  \item Noch wenig Traffic
  \item Ein Produkt, ein enger fachlicher Zusammenhang
\end{itemize}

Bearbeiten Sie die Fragen:
\begin{enumerate}
  \item Welche Architektur wäre als Startpunkt naheliegend?
  \item Nennen Sie \textbf{zwei Gründe}, warum diese Wahl sinnvoll ist.
  \item Nennen Sie \textbf{ein Risiko}, das bei späterem Wachstum entstehen könnte.
  \item Was müsste man tun, damit die gewählte Architektur später nicht zum Problem wird?
\end{enumerate}

\section*{4. Microservices kritisch beurteilen}
Viele Firmen sprechen von Microservices. Beantworten Sie die Fragen stichwortartig oder in kurzen Sätzen.

\begin{enumerate}
  \item Nennen Sie \textbf{zwei Vorteile} von Microservices.
  \item Nennen Sie \textbf{drei zusätzliche Herausforderungen} gegenüber einem Monolithen.
  \item Warum lösen Microservices Team- oder Organisationsprobleme nicht automatisch?
  \item In welchem Fall wären Microservices eher übertrieben?
\end{enumerate}

\section*{5. Client-Server vs. Peer-to-Peer}
Vergleichen Sie die beiden Kommunikationsmodelle.

\begin{tabularx}{\linewidth}{|>{\raggedright\arraybackslash}p{0.25\linewidth}|>{\raggedright\arraybackslash}X|>{\raggedright\arraybackslash}X|}
\hline
\textbf{Frage} & \textbf{Client-Server} & \textbf{Peer-to-Peer} \\
\hline
Wer stellt typischerweise die zentrale Logik bereit? & & \\[1.1cm] \hline
Wie erfolgt die Datenkontrolle oder Datenspeicherung? & & \\[1.1cm] \hline
Wo ist der Betrieb meist einfacher? & & \\[1.1cm] \hline
Nennen Sie ein typisches Praxisbeispiel. & & \\[1.1cm] \hline
\end{tabularx}

\section*{6. Architektur eines Online-Shops}
Ein Online-Shop benötigt folgende Funktionen:
\begin{itemize}
  \item Produktanzeige
  \item Warenkorb
  \item Checkout
  \item E-Mail-Benachrichtigungen
  \item Suchfunktion
\end{itemize}

Bearbeiten Sie dazu:
\begin{enumerate}
  \item Beschreiben Sie kurz, wie dieser Shop als \term{Monolith} aufgebaut sein könnte.
  \item Beschreiben Sie kurz, wie derselbe Shop als \term{Microservice-System} aufgebaut sein könnte.
  \item Welche Variante würden Sie für ein kleines Team zu Beginn bevorzugen?
  \item Welche Variante könnte bei starkem Wachstum interessanter werden und warum?
\end{enumerate}

\section*{7. Kombinationsverständnis}
Beantworten Sie die folgenden Aussagen mit \textbf{richtig / falsch} und begründen Sie kurz.

\begin{enumerate}
  \item Ein Monolith kann im Client-Server-Modell betrieben werden.
  \item Microservices bedeuten automatisch Peer-to-Peer.
  \item Peer-to-Peer-Systeme haben nie zentrale Elemente.
  \item Ein gut strukturierter Monolith kann für längere Zeit eine sinnvolle Lösung sein.
\end{enumerate}

\section*{8. Entwurfsaufgabe}
Sie planen eine Plattform für eine Schule mit folgenden Funktionen:
\begin{itemize}
  \item Lernmaterialien anzeigen
  \item Aufgaben hochladen
  \item Nachrichten zwischen Lehrpersonen und Lernenden
  \item Videokonferenzen einbinden
\end{itemize}

Bearbeiten Sie die Aufgaben:
\begin{enumerate}
  \item Zeichnen oder beschreiben Sie eine einfache Zielarchitektur.
  \item Entscheiden Sie sich begründet für einen \term{Monolithen} oder für eine erste \term{Microservice}-Aufteilung.
  \item Beschreiben Sie, wo in Ihrem Entwurf ein \term{Client-Server}-Muster vorkommt.
  \item Nennen Sie einen Teil, der eventuell später ausgelagert werden könnte.
  \item Begründen Sie, warum \term{Peer-to-Peer} hier eher passend oder eher unpassend wäre.
\end{enumerate}

\section*{Abgabe / Besprechung}
Die Aufgaben werden in der Lektion besprochen. Halten Sie Ihre Antworten so fest, dass Sie diese fachlich sauber vorstellen und begründen können.

\end{document}
